% Generated draft: requires mathematical and source review.
\documentclass[11pt]{article}
\usepackage[margin=25mm]{geometry}
\usepackage{fontspec}
\setmainfont{Latin Modern Roman}
\usepackage{amsmath,amssymb,mathtools,unicode-math}
\setmathfont{Latin Modern Math}
\usepackage{enumitem,needspace,xurl,hyperref}
\setlength{\parindent}{0pt}
\setlength{\parskip}{5pt}
\setlength{\emergencystretch}{3em}
\newcommand{\N}{\mathbb{N}}
\newcommand{\Z}{\mathbb{Z}}
\newcommand{\Q}{\mathbb{Q}}
\newcommand{\R}{\mathbb{R}}
\newcommand{\C}{\mathbb{C}}
\newcommand{\F}{\mathbb{F}}
\newcommand{\A}{\mathbb{A}}
\newcommand{\PP}{\mathbb{P}}
\newcommand{\E}{\mathbb{E}}
\newcommand{\eps}{\varepsilon}
\newcommand{\dd}{\,\mathrm{d}}
\begin{document}
\section*{UnsolvedMath ID: 10400127}
\textbf{Source problem code:} catalogue entry
\subsection{Definitions and mathematical statement}
For a closed three-manifold \(M\), let \(\tau_N^{SU(2)}(M)\) denote the source's \(SU(2)\) quantum invariant at the integer parameter \(N\). Let \(v_3\) be the hyperbolic volume of a regular ideal tetrahedron, and define the source's extended volume by
\[\operatorname{vol}(M)=v_3\lVert M\rVert,\]
where \(\lVert M\rVert\) is the simplicial volume. For an oriented closed manifold, this is the infimum of \(\sum_s|a_s|\) over real singular cycles \(\sum_s a_s\sigma_s\) representing its fundamental class. For hyperbolic \(M\), the source uses its hyperbolic volume; for nonhyperbolic \(M\), its footnote explicitly prescribes the simplicial-volume expression above.

The conjecture attributed to H. Murakami [294] is
\[2\pi i\;\operatorname{o\!\!\!\!-lim}_{N\to\infty}
\frac{\log\tau_N^{SU(2)}(M)}{N}
=\operatorname{CS}(M)+i\operatorname{vol}(M)\]
for every closed three-manifold \(M\), where the source calls \(\operatorname{o\!\!\!\!-lim}\) the \emph{optimistic limit} introduced in [294]. This is not a replacement by an ordinary limit. Its definition, logarithm/phase choices and any congruence convention in the displayed equality must be supplied from that reference. The normalization factor \(2\pi i\) is retained.

The footnote expressly makes the general-manifold statement conditional on an appropriate extension of \(\operatorname{CS}(M)\): it says that this invariant is defined there for hyperbolic manifolds and that a definition for arbitrary closed three-manifolds is itself conjectural, referring to Problem 7.16. Thus the formula for all \(M\) is a coupled extension-and-asymptotics conjecture, not an assertion that the required \(\operatorname{CS}(M)\) has already been specified. Orientation conventions are not stated in the short source assertion and must be fixed for the Chern--Simons term.
\subsection{Short English statement}
Murakami predicts that the optimistic large-parameter behavior of the \(SU(2)\) quantum invariant encodes Chern--Simons invariant plus imaginary volume. The volume extends through simplicial volume, while the needed Chern--Simons extension is itself part of the conjectural background.
\subsection{Sources}
\begin{itemize}
\item Local UnsolvedMath registry, permanent ID \(10400127\), exact source code \texttt{catalogue entry}.
\item H. Murakami [294], attribution in Ohtsuki, \emph{Problems on invariants of knots and 3-manifolds} (2002), Conjecture 7.12, PDF page 109, footnotes 23--24 and reference to Problem 7.16: \url{https://msp.org/gtm/2002/04/gtm-2002-04-024s.pdf}. The full entry for [294] is absent. Registry access 2026-07-29, pdf-font-structure extraction.
\end{itemize}
\end{document}
